\documentclass[12pt,a4paper]{article}
\usepackage[utf8]{inputenc}
\usepackage[spanish]{babel}
\usepackage{amsmath}
\usepackage{amsfonts}
\usepackage{amssymb}
\usepackage{geometry}
\usepackage{graphicx}
\usepackage{titlesec}
\usepackage{parskip}
\usepackage{fancyhdr}
\usepackage{enumitem}

\geometry{margin=1in}

% Configuración de encabezado y pie de página
\pagestyle{fancy}
\fancyhf{}
\fancyhead[L]{\textbf{Examen de Lectura}}
\fancyhead[C]{\textbf{Control 2}}
\fancyhead[R]{\today}
\fancyfoot[C]{\thepage}

% Formato de títulos
\titleformat{\section}
{\normalfont\Large\bfseries}{\thesection}{1em}{}
\titleformat{\subsection}
{\normalfont\large\bfseries}{\thesubsection}{1em}{}

% Título principal
\title{\vspace{-1cm}
\huge\textbf{Control de Lectura N°2} \\
\Large\textbf{Orientation and Decision-Making for Soccer Based on Sports Analytics and AI: A Systematic Review} \\
\vspace{0.5cm}
\large{Pu, Pan, Wang, Liu, Chen, Ma \& Cui (2024)}}
\author{\textbf{Armando Arredondo Valle} \\
\textit{Magíster en Inteligencia Artificial} \\
\textit{Pontificia Universidad Católica de Chile}}
\date{30 de abril de 2026}

\begin{document}

\maketitle
\thispagestyle{fancy}

\vspace{0.5cm}

\vspace{0.5cm}
\hrule
\vspace{0.5cm}

\section*{Pregunta N°1}
\textbf{La lectura plantea que el análisis del fútbol puede entenderse mediante un ciclo de observación, orientación, decisión y acción. Explique brevemente qué significa este ciclo en el contexto del fútbol y cómo se relaciona con la Ciencia de Datos y la Inteligencia Artificial.}

\subsection*{Respuesta}
A lo largo de la lectura se plantea justamente que el análisis del fútbol puede modelarse mediante el ciclo OODA (Observation-Orientation-Decision-Action), un modelo originalmente creado para pilotos de combate aéreo, pero que se adapta de forma interesante al contexto deportivo.\\

En el fútbol, la \textit{observación} corresponde a la recolección de datos mediante sistemas de tracking óptico, cámaras y dispositivos GPS. La \textit{orientación} es donde entran los modelos de análisis como los expected goals (xG), el expected possession value (EPV) y el pitch control. La \textit{decisión} representa la etapa donde se generan recomendaciones tácticas, pudiendo incluso que en entornos virtuales como Google Research Football los algoritmos controlen agentes directamente. Finalmente, la \textit{acción} es la ejecución a través del entrenamiento diario y los partidos.\\

Es fascinante cómo este ciclo es prácticamente un pipeline de ciencia de datos aplicado al deporte: recolectar, analizar, decidir y actuar. De igual forma, el fútbol ha sido un benchmark representativo para la IA, especialmente en aprendizaje por refuerzo multiagente, donde la cooperación entre 11 jugadores y las recompensas escasas representan problemas de frontera.

\vspace{0.5cm}
\hrule
\vspace{0.5cm}

\section*{Pregunta N°2}
\textbf{Según la lectura, existen distintos tipos de datos usados en el análisis del fútbol, como event data, tracking data y otros datos complementarios. Compare estos tipos de datos e indique qué ventajas y limitaciones tiene cada uno para construir modelos de análisis deportivo.}

\subsection*{Respuesta}
La lectura nos habla sobre tres tipos principales de datos en el análisis del fútbol.\\

Los \textbf{event data} son un registro cronológico de eventos clave: pases, tiros, faltas y tackles. Proveedores como Opta, WyScout y StatsBomb ofrecen estos datos, incluso con datasets abiertos de casi 2000 partidos. Su ventaja es la disponibilidad y volumen, permitiendo calcular métricas como xG y expected threat (xT). Sin embargo, el gran problema aquí es que no capturan información de jugadores fuera del balón, y los intervalos entre eventos son de varios segundos.\\

Los \textbf{tracking data} registran la posición de todos los jugadores y el balón a tasa fija (por ejemplo, 10 fps), permitiendo obtener velocidad, aceleración y tácticas de jugadores sin balón. Con estos datos se construyen modelos como el de pitch control. A fin de cuentas, su limitación principal es el costo; los datasets abiertos son muy escasos.\\

Los \textbf{otros datos} incluyen información física (biométrica, sprints) y datos de redes sociales. Permiten evaluar rendimiento físico individual pero su alcance es más limitado.\\

De igual forma, la lectura muestra que la combinación de event data con tracking data es donde se obtienen los mejores resultados, pudiendo tener análisis tácticos más completos.

\vspace{0.5cm}
\hrule
\vspace{0.5cm}

\section*{Pregunta N°3}
\textbf{Imagine que usted trabaja con datos del fútbol chileno. Proponga una visualización que permita comunicar una idea relevante inspirada en la lectura. Indique qué datos usaría, qué gráfico construiría y qué pregunta deportiva ayudaría a responder.}

\subsection*{Respuesta}
Inspirándome en el modelo de pitch control de la lectura, propondría un \textbf{mapa de calor de control territorial} para equipos de la liga chilena.\\

Usaría tracking data de los partidos (posiciones de jugadores y balón), complementado con event data de pases, tiros y recuperaciones. El gráfico sería un mapa de calor sobre la cancha donde los colores representen la ventaja territorial: rojo donde el equipo local domina, azul donde domina el visitante y tonos neutros en zonas de equilibrio, utilizando el modelo de influencia individual con distribuciones normales que describe la lectura.\\

La pregunta deportiva sería: \textit{¿Cómo varía el control territorial de un equipo chileno jugando de local versus visitante, y qué zonas son las más disputadas en los clásicos?} Siendo así que esta visualización permitiría a los entrenadores identificar zonas débiles del rival y momentos críticos donde se pierde el control espacial.

\vspace{0.5cm}
\hrule
\vspace{0.5cm}

\section*{Pregunta N°4}
\textbf{La lectura muestra cómo la IA y la Ciencia de Datos pueden apoyar decisiones en un contexto complejo como el fútbol. Explique cómo una idea de la lectura podría aplicarse a su trabajo diario, aunque no esté relacionado directamente con deporte.}

\subsection*{Respuesta}
A lo largo de la lectura se presenta un concepto aplicable a mi trabajo como ingeniero de software en Microsoft: el ciclo OODA y los modelos integrados de conocimiento y datos.\\

El ciclo OODA del fútbol se traduce casi directo a lo que hago en operaciones live-site en Azure. La \textit{observación} es nuestra capa de observabilidad con Azure Monitor y Geneva, como el tracking captura posiciones de jugadores. La \textit{orientación} es el análisis de esos datos para detectar anomalías, similar a cómo el pitch control evalúa la ventaja espacial. La \textit{decisión} es el diagnóstico del incidente, donde al igual que en el fútbol, a fin de cuentas el juicio del ingeniero es insustituible. Y la \textit{acción} es el despliegue de correcciones a través de nuestros pipelines CI/CD con staged deployments y slot swaps.\\

De igual forma, el análisis contrafactual de la lectura me recuerda a los postmortems de incidentes: analizamos qué habría ocurrido con una validation gate adicional o una estrategia de disaster recovery diferente.\\

A fin de cuentas, la lección más valiosa es que la combinación de conocimiento experto con técnicas de IA produce resultados que ninguno lograría por separado, y eso se refleja claramente en cómo la experiencia del ingeniero combinada con alertas inteligentes nos permite mantener servicios distribuidos a escala.

\end{document}